\documentclass[10pt,a4paper]{article}

% =========================================================
% PACKAGES
% =========================================================

\usepackage[utf8]{inputenc}
\usepackage[T1]{fontenc}
\usepackage[english]{babel}

\usepackage[
left=25mm,
right=25mm,
top=20mm,
bottom=22mm,
headheight=14pt,
headsep=10mm,
footskip=10mm
]{geometry}

\usepackage{lmodern}
\usepackage{microtype}

\usepackage{xcolor}
\usepackage{graphicx}
\usepackage{booktabs}
\usepackage{array}
\usepackage{enumitem}
\usepackage{listings}
\usepackage{fancyhdr}
\usepackage{titlesec}
\usepackage{tcolorbox}
\usepackage{hyperref}
\usepackage{setspace}

% =========================================================
% COLORS
% =========================================================

\definecolor{THGABlue}{RGB}{35,79,120}
\definecolor{THGADarkBlue}{RGB}{15,65,110}
\definecolor{THGARed}{RGB}{205,0,0}
\definecolor{THGAGreen}{RGB}{75,180,160}
\definecolor{THGABoxBlue}{RGB}{205,225,245}
\definecolor{CodeBackground}{RGB}{225,225,225}
\definecolor{CodeBlue}{RGB}{0,70,130}
\definecolor{CodeGreen}{RGB}{0,110,60}
\definecolor{CodeRed}{RGB}{150,0,0}

% =========================================================
% PAGE STYLE
% =========================================================

\pagestyle{fancy}

\fancyhf{}

\fancyhead[L]{%
	\textcolor{THGABlue}{\small\textbf{Databases}}
}

\fancyhead[R]{%
	\textcolor{gray}{\small
		Stephan Bökelmann $\cdot$ Summer Term 2026}
}

\fancyfoot[L]{%
	\textcolor{THGARed}{\small sboekelmann@ep1.rub.de}
}

\fancyfoot[C]{%
	\textcolor{gray}{\small THGA Bochum}
}

\fancyfoot[R]{%
	\textcolor{gray}{\small\thepage}
}

\renewcommand{\headrulewidth}{0.8pt}
\renewcommand{\headrule}{%
	\hbox to\headwidth{%
		\color{THGABlue}\leaders\hrule height \headrulewidth\hfill
	}%
}

\renewcommand{\footrulewidth}{0pt}

% =========================================================
% SECTION STYLE
% =========================================================

\titleformat{\section}
{\large\bfseries\color{THGABlue}}
{\thesection}
{0.7em}
{}
[\vspace{1mm}\color{THGABlue}\titlerule]

\titleformat{\subsection}
{\normalsize\bfseries\color{THGABlue}}
{\thesubsection}
{0.6em}
{}

\titlespacing*{\section}
{0pt}{10pt}{5pt}

\titlespacing*{\subsection}
{0pt}{7pt}{4pt}

% =========================================================
% HYPERLINKS
% =========================================================

\hypersetup{
	colorlinks=true,
	linkcolor=THGABlue,
	urlcolor=THGARed,
	citecolor=THGABlue
}

% =========================================================
% LISTINGS
% =========================================================

\lstset{
	basicstyle=\ttfamily\scriptsize,
	backgroundcolor=\color{CodeBackground},
	frame=none,
	numbers=left,
	numberstyle=\tiny\color{gray},
	numbersep=7pt,
	breaklines=true,
	breakatwhitespace=false,
	showstringspaces=false,
	keepspaces=true,
	columns=fullflexible,
	tabsize=4
}

% =========================================================
% BOXES
% =========================================================

\newtcolorbox{infobox}{
	colback=THGABoxBlue,
	colframe=THGABlue,
	boxrule=1.2pt,
	arc=0pt,
	left=5pt,
	right=5pt,
	top=5pt,
	bottom=5pt
}

\newtcolorbox{greenbox}{
	colback=white,
	colframe=THGAGreen,
	boxrule=1.2pt,
	arc=0pt,
	left=5pt,
	right=5pt,
	top=5pt,
	bottom=5pt
}

% =========================================================
% DOCUMENT INFORMATION
% =========================================================

\title{
	\textcolor{THGABlue}{
		\textbf{Term Project -- Documentation}
	}
}

\author{}
\date{}

% =========================================================
% DOCUMENT
% =========================================================

\begin{document}
	
	% =========================================================
	% TITLE
	% =========================================================
	
	\thispagestyle{empty}
	
	% Dark blue title bar
	\begin{tcolorbox}[
		colback=THGADarkBlue,
		colframe=THGADarkBlue,
		boxrule=0pt,
		arc=0pt,
		left=8pt,
		right=8pt,
		top=8pt,
		bottom=8pt
		]
		\centering
		{\Large\color{white}\textbf{Term Project -- Documentation}}
	\end{tcolorbox}
	
	\vspace{7mm}
	
	{\large
		\textcolor{THGABlue}{
			\textbf{Study Room Reservation System: a database-backed reservation system}
		}
	}
	
	\vspace{6mm}
	
	\textcolor{THGARed}{\rule{\textwidth}{1.5pt}}
	
	\vspace{6mm}
	
	\noindent
	Dozent: Stephan Bökelmann
	\hspace{5mm}
	\textcolor{THGARed}{sboekelmann@ep1.rub.de}
	\hspace{5mm}
	Technische Hochschule Georg Agricola, Bochum
	
	\vspace{7mm}
	
	\begin{infobox}
		
		\textbf{Author:} Imane Tamouh
		
		\textbf{Module:} Introduction to Database Management Systems
		$\cdot$ THGA Bochum
		
		\textbf{Stack:} PostgreSQL $\cdot$ FastAPI $\cdot$
		Docker Compose $\cdot$ tkinter (.deb)
		
		\textbf{Repository:}
		\textcolor{THGARed}{
			\texttt{https://github.com//imanetamouh7-bit/study-room-reservation>}
		}
		
	\end{infobox}
	
	\vspace{7mm}
	
	% =========================================================
	% TABLE OF CONTENTS
	% =========================================================
	
	{\large\bfseries\textcolor{THGABlue}{Inhaltsverzeichnis}}
	
	\vspace{1mm}
	
	\textcolor{THGABlue}{\rule{\textwidth}{0.7pt}}
	
	\vspace{2mm}
	
	\tableofcontents
	
	\newpage
	
	
	% =========================================================
	% 1 INTRODUCTION
	% =========================================================
	
	\section{Introduction and motivation}
	
	Students often need a quiet place to study, work on assignments or
	prepare for examinations. If study rooms are managed manually, it
	can be difficult to determine which rooms are available and whether
	a room has already been reserved for a particular period.
	
	The Study Room Reservation System replaces manual reservation
	management with a relational database system. It manages students,
	study rooms, equipment and reservations, exposes the data through a
	REST API and provides a small desktop frontend.
	
	\begin{greenbox}
		
		\textbf{Usage scenario.}
		A student wants to reserve a study room for an examination
		preparation session. The student opens the frontend, enters the API
		URL and the X-API-Key and connects to the backend.
		
		The available study rooms are displayed together with their
		capacity. The student selects a room, enters the date and time and
		creates the reservation.
		
		If the selected room is already reserved for the requested time
		period, the system rejects the reservation and informs the student
		that another room or another time period has to be selected.
		
	\end{greenbox}
	
	
	% =========================================================
	% 2 DATA MODEL
	% =========================================================
	
	\section{Data model}
	
	The domain has five entity types: students, study rooms, equipment,
	reservations and the associative entity connecting rooms with
	equipment.
	
	A student can make many reservations over time, while each
	reservation belongs to one student. A room can be reserved many
	times over different time periods, while each reservation refers to
	one room.
	
	A room can have many equipment types and an equipment type can occur
	in many rooms. This is an N:M relationship that is resolved by the
	\texttt{room\_equipment} relation.
	
	\begin{greenbox}
		
		\textbf{ER model (textual).}
		
		\begin{itemize}[leftmargin=5mm,itemsep=1pt]
			\item \textit{student} (1) -- makes -- (N) \textit{reservation}
			\item \textit{room} (1) -- is reserved in -- (N) \textit{reservation}
			\item \textit{room} (1) -- is assigned -- (N) \textit{room\_equipment}
			\item \textit{equipment} (1) -- occurs in -- (N) \textit{room\_equipment}
		\end{itemize}
		
		Thus \textit{room} N:M \textit{equipment}, resolved through
		\textit{room\_equipment}.
		
	\end{greenbox}
	
	
	% =========================================================
	% 2.1 RELATIONAL SCHEMA
	% =========================================================
	
	\subsection{Relational schema}
	
	Derived from the ER model, with primary keys underlined and foreign
	keys in italics:
	
	\begin{itemize}[leftmargin=6mm,itemsep=3pt]
		
		\item
		\textit{student}(
		\underline{student\_id},
		name,
		email
		)
		
		\item
		\textit{room}(
		\underline{room\_id},
		room\_number,
		capacity
		)
		
		\item
		\textit{equipment}(
		\underline{equipment\_id},
		name
		)
		
		\item
		\textit{reservation}(
		\underline{reservation\_id},
		date,
		start\_time,
		end\_time,
		\textit{student\_id},
		\textit{room\_id}
		)
		
		\item
		\textit{room\_equipment}(
		\underline{\textit{room\_id}, \textit{equipment\_id}}
		)
		
	\end{itemize}
	
	\textbf{Normalisation.}
	Every non-key attribute depends only on the key of its table, and
	there are no transitive dependencies between non-key attributes.
	
	Student information is stored only in \textit{student}, room
	information only in \textit{room}, and equipment information only in
	\textit{equipment}. Reservations reference students and rooms by
	their keys instead of duplicating their data.
	
	The N:M relationship between rooms and equipment is represented by
	the separate \textit{room\_equipment} relation.
	
	The schema is therefore in 3NF.
	
	
	% =========================================================
	% 3 DATABASE
	% =========================================================
	
	\section{Database (PostgreSQL)}
	
	The schema and seed data are created by SQL init scripts that
	PostgreSQL runs when the database is initially created.
	
	The PostgreSQL service uses version 16 and stores its data in the
	Docker volume \texttt{pg\_data}.
	
	\begin{lstlisting}[language=SQL,
		caption={schema.sql (excerpt) -- student, room and equipment}]
		CREATE TABLE student (
		student_id INTEGER GENERATED ALWAYS AS IDENTITY PRIMARY KEY,
		name VARCHAR(120) NOT NULL,
		email VARCHAR(200) NOT NULL UNIQUE
		);
		
		CREATE TABLE room (
		room_id INTEGER GENERATED ALWAYS AS IDENTITY PRIMARY KEY,
		room_number VARCHAR(20) NOT NULL UNIQUE,
		capacity INTEGER NOT NULL CHECK (capacity > 0)
		);
		
		CREATE TABLE equipment (
		equipment_id INTEGER GENERATED ALWAYS AS IDENTITY PRIMARY KEY,
		name VARCHAR(100) NOT NULL UNIQUE
		);
	\end{lstlisting}
	
	
	\begin{lstlisting}[language=SQL,
		caption={schema.sql (excerpt) -- reservation}]
		CREATE TABLE reservation (
		reservation_id INTEGER GENERATED ALWAYS AS IDENTITY PRIMARY KEY,
		date DATE NOT NULL,
		start_time TIME NOT NULL,
		end_time TIME NOT NULL,
		student_id INTEGER NOT NULL,
		room_id INTEGER NOT NULL,
		
		CHECK (end_time > start_time),
		
		FOREIGN KEY (student_id)
		REFERENCES student (student_id)
		ON DELETE RESTRICT
		ON UPDATE CASCADE,
		
		FOREIGN KEY (room_id)
		REFERENCES room (room_id)
		ON DELETE RESTRICT
		ON UPDATE CASCADE
		);
	\end{lstlisting}
	
	
	\begin{lstlisting}[language=SQL,
		caption={schema.sql (excerpt) -- room equipment}]
		CREATE TABLE room_equipment (
		room_id INTEGER NOT NULL,
		equipment_id INTEGER NOT NULL,
		
		PRIMARY KEY (room_id, equipment_id),
		
		FOREIGN KEY (room_id)
		REFERENCES room (room_id)
		ON DELETE CASCADE
		ON UPDATE CASCADE,
		
		FOREIGN KEY (equipment_id)
		REFERENCES equipment (equipment_id)
		ON DELETE CASCADE
		ON UPDATE CASCADE
		);
	\end{lstlisting}
	
	The \texttt{student} table stores four initial students. The
	\texttt{room} table contains four study rooms with capacities 4, 6,
	8 and 12.
	
	The equipment table contains Projector, Whiteboard, Monitor and
	Video Conference.
	
	The seed data also creates room-equipment assignments and four
	initial reservations.
	
	
	% =========================================================
	% 4 REST API
	% =========================================================
	
	\section{REST API (FastAPI)}
	
	The API is the only component that talks to the database. Read
	endpoints retrieve information from PostgreSQL, while the write
	endpoint requires the header \texttt{X-API-Key}.
	
	\begin{center}
		\begin{tabular}{llll}
			\toprule
			\textbf{Method} & \textbf{Path} & \textbf{Purpose} & \textbf{Key}\\
			\midrule
			GET  & \texttt{/} &
			API status &
			-- \\
			
			GET  & \texttt{/rooms} &
			list study rooms &
			-- \\
			
			GET  & \texttt{/reservations} &
			reservations + JOIN &
			-- \\
			
			POST & \texttt{/reservations} &
			create reservation &
			yes \\
			
			GET & \texttt{/rooms/usage} &
			room usage (JOIN + agg.) &
			-- \\
			
			\bottomrule
		\end{tabular}
	\end{center}
	
	The \texttt{/reservations} endpoint demonstrates a JOIN between
	reservation, student and room.
	
	The \texttt{/rooms/usage} endpoint demonstrates a JOIN combined
	with an aggregation. It counts the reservations belonging to each
	room.
	
	\begin{lstlisting}[language=SQL,
		caption={Query behind GET /rooms/usage (JOIN + aggregation)}]
		SELECT
		ro.room_id,
		ro.room_number,
		ro.capacity,
		COUNT(r.reservation_id) AS reservation_count
		FROM room AS ro
		LEFT JOIN reservation AS r
		ON ro.room_id = r.room_id
		GROUP BY
		ro.room_id,
		ro.room_number,
		ro.capacity
		ORDER BY ro.room_id;
	\end{lstlisting}
	
	The write path is guarded by a dependency that checks the API key.
	
	\begin{lstlisting}[language=Python,
		caption={API-key guard}]
		def verify_api_key(x_api_key: str = Header(...)):
		if x_api_key != ENV["API_KEY"]:
		raise HTTPException(
		status_code=401,
		detail="Invalid API key",
		)
	\end{lstlisting}
	
	The reservation creation endpoint uses this dependency:
	
	\begin{lstlisting}[language=Python,
		caption={Protected POST /reservations}]
		@app.post("/reservations", status_code=201)
		def create_reservation(
		reservation: ReservationCreate,
		_: None = Depends(verify_api_key),
		):
		...
	\end{lstlisting}
	
	Before inserting a reservation, the backend checks whether another
	reservation for the same room and date overlaps with the requested
	time interval.
	
	If an overlap exists, the API returns HTTP 409:
	
	\begin{lstlisting}[language=Python,
		caption={Reservation conflict}]
		raise HTTPException(
		status_code=409,
		detail="Room is already reserved for this time period",
		)
	\end{lstlisting}
	
	
	% =========================================================
	% 5 FRONTEND
	% =========================================================
	
	\section{Frontend (tkinter, packaged as .deb)}
	
	The frontend is a small tkinter application with screens for
	available study rooms, new reservations, existing reservations and
	room usage.
	
	It never touches the database directly -- it only calls the API.
	
	At startup a connection dialog asks for the API URL and the
	X-API-Key.
	
	The current default API URL is:
	
	\begin{lstlisting}
		http://127.0.0.1:8000
	\end{lstlisting}
	
	After a successful connection, the frontend displays the available
	study rooms.
	
	The main screen provides four actions:
	
	\begin{itemize}[leftmargin=6mm,itemsep=2pt]
		\item New Reservation
		\item View Reservations
		\item Room Usage
		\item Disconnect
	\end{itemize}
	
	\subsection*{New Reservation}
	
	The reservation form contains:
	
	\begin{itemize}[leftmargin=6mm,itemsep=2pt]
		\item Date
		\item Start time
		\item End time
		\item Student ID
		\item Study room
	\end{itemize}
	
	The frontend sends the values to:
	
	\begin{lstlisting}
		POST /reservations
	\end{lstlisting}
	
	A successful request displays the generated reservation ID.
	
	If the backend returns HTTP 409, the frontend displays a warning
	that the room is already reserved for the selected time period.
	
	\subsection*{View Reservations}
	
	The reservation view calls:
	
	\begin{lstlisting}
		GET /reservations
	\end{lstlisting}
	
	The returned data is displayed in a table containing:
	
	\begin{itemize}[leftmargin=6mm,itemsep=2pt]
		\item reservation ID
		\item date
		\item start time
		\item end time
		\item student
		\item room
	\end{itemize}
	
	\subsection*{Room Usage}
	
	The room usage view calls:
	
	\begin{lstlisting}
		GET /rooms/usage
	\end{lstlisting}
	
	The table displays the room ID, room number, capacity and number of
	reservations.
	
	\subsection*{Packaging}
	
	The frontend source contains a PyInstaller specification file
	\texttt{main.spec}. The current project also contains generated
	frontend files in the \texttt{build} and \texttt{dist} directories.
	
	For the final deployment, the frontend is packaged as a Debian
	package (\texttt{.deb}) so that it can be installed independently
	from the Docker Compose backend.
	
	
	% =========================================================
	% 6 OPERATION AND DEPLOYMENT
	% =========================================================
	
	\section{Operation and deployment}
	
	The backend runs as two containers, orchestrated by Docker Compose:
	a PostgreSQL database container and a FastAPI API container.
	
	Credentials and configuration values are supplied through the
	\texttt{.env} file.
	
	\begin{lstlisting}[caption={docker-compose.yml (excerpt)}]
		services:
		
		postgres:
		
		image: postgres:16
		
		restart: unless-stopped
		
		environment:
		POSTGRES_USER: ${POSTGRES_USER}
		POSTGRES_PASSWORD: ${POSTGRES_PASSWORD}
		POSTGRES_DB: ${POSTGRES_DB}
		
		volumes:
		- pg_data:/var/lib/postgresql/data
		- ./database/schema.sql:/docker-entrypoint-initdb.d/01-schema.sql
		- ./database/seed.sql:/docker-entrypoint-initdb.d/02-seed.sql
		
		healthcheck:
		test:
		[
		"CMD-SHELL",
		"pg_isready -U ${POSTGRES_USER} -d ${POSTGRES_DB}"
		]
		interval: 5s
		timeout: 5s
		retries: 5
		
		api:
		
		build: ./backend
		
		restart: unless-stopped
		
		env_file:
		- .env
		
		environment:
		DB_HOST: postgres
		DB_NAME: ${POSTGRES_DB}
		DB_USER: ${POSTGRES_USER}
		DB_PASSWORD: ${POSTGRES_PASSWORD}
		API_KEY: ${API_KEY}
		
		ports:
		- "8000:8000"
		
		depends_on:
		postgres:
		condition: service_healthy
		
		volumes:
		pg_data:
	\end{lstlisting}
	
	The PostgreSQL container uses the named volume
	\texttt{pg\_data} to persist the database data.
	
	The API container is built from the \texttt{backend} directory and
	is started only after the PostgreSQL health check succeeds.
	
	\subsection*{Deployment steps}
	
	\begin{enumerate}[leftmargin=7mm,itemsep=3pt]
		
		\item
		Start the database and API:
		
		\begin{lstlisting}
			docker compose up -d
		\end{lstlisting}
		
		\item
		Check the running containers:
		
		\begin{lstlisting}
			docker compose ps
		\end{lstlisting}
		
		\item
		Install the frontend Debian package:
		
		\begin{lstlisting}
			sudo dpkg -i study-room-reservation_0.1.0_amd64.deb
		\end{lstlisting}
		
		\item
		Launch the frontend.
		
		\item
		Enter the API URL and the X-API-Key in the connection dialog.
		
	\end{enumerate}
	
	The FastAPI documentation can be accessed locally at:
	
	\begin{lstlisting}
		http://127.0.0.1:8000/docs
	\end{lstlisting}
	
	
	% =========================================================
	% 7 REFLECTION
	% =========================================================
	
	\section{Reflection}
	
	The separation into database, API and frontend provides a clear
	architecture for the Study Room Reservation System.
	
	The database stores the application data and enforces important
	integrity constraints. The API provides the interface between the
	database and the graphical application.
	
	The JOIN query used by the reservation endpoint demonstrates how
	related information can be retrieved without duplicating student and
	room data. The aggregation query used by the room usage endpoint
	demonstrates how reservation counts can be calculated directly in
	SQL.
	
	The reservation conflict check was an important part of the
	application because two overlapping reservations must not be
	created for the same room.
	
	Docker Compose simplifies the operation of the backend by starting
	PostgreSQL and FastAPI together. The health check ensures that the
	API waits for the database to become available.
	
	The tkinter frontend provides the user with a simple interface for
	connecting to the API, viewing rooms and reservations, creating new
	reservations and viewing room usage.
	
	
	% =========================================================
	% SOURCES AND REUSE
	% =========================================================
	
	\section*{Sources and reuse}
	
	Reused with attribution, as permitted:
	
	\begin{itemize}[leftmargin=6mm,itemsep=3pt]
		
		\item
		API-key guard and REST API structure -- adapted from the
		DBMS course examples and DBMS project exercises.
		
		\item
		Docker Compose structure with PostgreSQL and FastAPI -- based
		on the DBMS Docker and FastAPI lectures.
		
		\item
		tkinter connection dialog and Treeview-based frontend layout --
		based on the frontend concepts presented in the DBMS lectures.
		
		\item
		PyInstaller and Debian packaging approach -- based on the
		course requirements for the frontend deployment.
		
	\end{itemize}
	
	The database schema, seed data and Study Room Reservation System
	application-specific functionality were developed for this project.
	
\end{document}